\documentclass[a4paper,11pt]{article}
\usepackage[a4paper,margin=2cm]{geometry}
\usepackage[T1]{fontenc}
\usepackage[utf8]{inputenc}
\usepackage[ngerman,english]{babel}
\usepackage{lmodern}
\usepackage{microtype}
\usepackage{xcolor}
\usepackage{parskip}
\usepackage{enumitem}
\usepackage[most]{tcolorbox}
\usepackage{titlesec}
\usepackage[hidelinks]{hyperref}
\usepackage{array}
\usepackage{longtable}
\usepackage[table]{xcolor}
\definecolor{HeaderBlueDark}{RGB}{70,110,170}
\definecolor{RowBlueLight}{RGB}{235,243,255}
\definecolor{RowBlueDark}{RGB}{220,233,250}
\definecolor{SoftGray}{RGB}{90,90,90}
\setlength{\parindent}{0pt}
\setlist[itemize]{leftmargin=1.4em,itemsep=0.35em,topsep=0.35em}
\linespread{1.06}
\renewcommand{\arraystretch}{1.35}
\setlength{\tabcolsep}{5pt}
\titleformat{\section}[block]{\Large\bfseries\color{HeaderBlueDark}}{}{0pt}{}
\titlespacing{\section}{0pt}{1.3em}{0.8em}
\titleformat{\subsection}[block]{\large\bfseries\color{HeaderBlueDark}}{}{0pt}{}
\titlespacing{\subsection}{0pt}{1.0em}{0.6em}
\newtcolorbox{exbox}{enhanced,breakable,colback=RowBlueLight,colframe=RowBlueDark,boxrule=0.4pt,arc=1.5mm,left=1.2mm,right=1.2mm,top=1mm,bottom=1mm,title={Beispiele \textnormal{(Examples)}},fonttitle=\bfseries,colbacktitle=RowBlueDark,coltitle=black}
\NewDocumentEnvironment{grammarentry}{mm+b}{%
\begin{tcolorbox}[enhanced,breakable,colback=white,colframe=HeaderBlueDark,boxrule=0.6pt,arc=1.5mm,left=1.5mm,right=1.5mm,top=1mm,bottom=1mm,colbacktitle=HeaderBlueDark,coltitle=white,fonttitle=\bfseries,title={#1\hfill\normalfont\footnotesize #2}]
#3
\end{tcolorbox}}{}
\newcommand{\GermanText}[1]{\textbf{Deutsch: }#1\par}
\newcommand{\EnglishText}[1]{{\small\color{SoftGray}\textbf{English: }#1\par}}
\newcommand{\ExampleItem}[2]{\item \textbf{#1}\\{\small\color{SoftGray}#2}}
\newcommand{\DEEN}[2]{\textbf{#1}\par{\footnotesize\color{SoftGray}#2}}
\title{Passiv mit Modalverben\\[0.3em]\large\textnormal{Passive with modal verbs}}
\date{}
\begin{document}
\maketitle
\tableofcontents
\newpage

\section{Grundbildung -- Basic formation}
\begin{grammarentry}{Modalverb + Partizip II + werden}{Modal verb + past participle + werden}
\GermanText{Im Hauptsatz steht das konjugierte Modalverb an Position 2. Am Satzende stehen \textbf{Partizip II + werden}.}
\EnglishText{In a main clause, the conjugated modal verb stands in position 2. At the end stand \textbf{past participle + werden}.}
\GermanText{Beispiel: \textbf{Die Aufgabe muss gemacht werden.}}
\EnglishText{Example: \textbf{The task must be done.}}
\end{grammarentry}

\section{Modalverben -- Modal verbs}
\begin{exbox}\begin{itemize}
\ExampleItem{Das Problem kann gelöst werden.}{The problem can be solved.}
\ExampleItem{Die Rechnung muss bezahlt werden.}{The bill must be paid.}
\ExampleItem{Die Tür darf nicht geöffnet werden.}{The door may not be opened.}
\ExampleItem{Der Antrag soll heute geprüft werden.}{The application is supposed to be checked today.}
\ExampleItem{Das Formular muss ausgefüllt werden.}{The form must be filled out.}
\end{itemize}\end{exbox}

\section{Wortstellung -- Word order}
\begin{grammarentry}{Hauptsatz}{Main clause}
\GermanText{\textbf{Die Aufgabe muss heute gemacht werden.} Das Modalverb ist konjugiert; \textbf{gemacht werden} steht am Ende.}
\EnglishText{\textbf{The task must be done today.} The modal verb is conjugated; \textbf{gemacht werden} stands at the end.}
\end{grammarentry}
\begin{grammarentry}{Nebensatz}{Subordinate clause}
\GermanText{Im Nebensatz steht das konjugierte Modalverb am Ende: \textbf{..., weil die Aufgabe heute gemacht werden muss.}}
\EnglishText{In a subordinate clause, the conjugated modal verb stands at the end: \textbf{..., because the task must be done today.}}
\end{grammarentry}

\section{Zeitformen -- Tenses}
\begin{grammarentry}{Präsens und Präteritum}{Present and simple past}
\GermanText{Präsens: \textbf{Die Aufgabe muss gemacht werden.} Präteritum: \textbf{Die Aufgabe musste gemacht werden.}}
\EnglishText{Present: \textbf{The task must be done.} Simple past: \textbf{The task had to be done.}}
\end{grammarentry}
\begin{grammarentry}{Perfekt -- fortgeschritten}{Perfect -- advanced}
\GermanText{Im Perfekt entsteht eine Ersatzinfinitiv-Konstruktion: \textbf{Die Aufgabe hat gemacht werden müssen.} Diese Form ist korrekt, aber für B1 nicht besonders wichtig und wird häufig vermieden.}
\EnglishText{In the perfect tense, a substitute-infinitive construction is formed: \textbf{The task had to be done.} This form is correct, but it is not especially important at B1 level and is often avoided.}
\end{grammarentry}

\section{Aktiv und Passiv -- Active and passive}
\begin{exbox}\begin{itemize}
\ExampleItem{Aktiv: Man muss die Rechnung bezahlen.}{Active: One must pay the bill.}
\ExampleItem{Passiv: Die Rechnung muss bezahlt werden.}{Passive: The bill must be paid.}
\end{itemize}\end{exbox}

\section{Merksatz -- Key rule}
\begin{grammarentry}{Kurz}{In short}
\GermanText{\textbf{Modalverb + Partizip II + werden}. Im Nebensatz kommt das konjugierte Modalverb ganz ans Ende.}
\EnglishText{\textbf{Modal verb + past participle + werden}. In a subordinate clause, the conjugated modal verb goes to the very end.}
\end{grammarentry}
\end{document}
